\documentclass[12pt]{article}
\usepackage[utf8]{inputenc}
\usepackage[T1]{fontenc}
\usepackage{hyperref}
\title{Wawonii Language Resources for K-12 Curriculum Development}
\author{Compiled by Claude Sonnet 5.5 (claude-sonnet-5-5) for Andr\'as Kornai}
\date{}
\begin{document}
\maketitle

\section{Introduction}
Wawonii (ISO 639-3: wow; Glottocode: wawo1239), also spelled Wowoni or Wowonii, is an Austronesian language of the Celebic branch. It is spoken on Wawonii Island in Southeast Sulawesi and on nearby Menui Island, with about 30,000 speakers. It is closely related to Bungku. A 2026 study describes it as threatened, with two dialects, Wawonii and Menui. A 2019 study reports that it is still used for everyday communication in Konawe Kepulauan regency, including by non-Wawonii residents.

\section{Primary Educational Resources}
I found no primers, textbooks or school programmes in search results.

\section{Linguistic and Reference Materials}
\begin{itemize}
\item ``Morphological Complexity in Wawonii: A Documentation of Affixation Patterns in a Threatened Language'', \emph{Theory and Practice in Language Studies}: \url{https://tpls.academypublication.com/index.php/tpls/article/view/12526}
\item Full text of the same article (PDF): \url{https://tpls.academypublication.com/index.php/tpls/article/download/12526/10483/40653}
\item Article in \emph{Ganaya} on use of the language in Konawe Kepulauan, 2019 (Jayapangus Press): \url{https://jayapanguspress.penerbit.org/index.php/ganaya/article/download/305/299/595}
\item Hasanuddin University publication record: \url{https://discover.unhas.ac.id/publications/pub_scopus_105035412367}
\item Wikipedia article on the Wawonii language: \url{https://en.wikipedia.org/wiki/Wawonii_language}
\item Omniglot page, noting that the language can be written in the Latin alphabet: \url{https://omniglot.com/writing/wawonii.htm}
\item Verbix page on Wawonii, which gives about 75\% lexical similarity with Bungku: \url{https://docs.verbix.com/Languages/Wawonii}
\end{itemize}

\section{Religious and Community Resources}
An Asia Harvest people-group profile exists: \url{https://people-groups.asiaharvest.org/Indonesia/Wawonii.pdf}. I did not check what it contains.

\section{Digital and Audio Resources}
No audio, video or corpus resources were found.

\section{Curriculum Development Context}
Wawonii has a reasonably large speaker base and recent academic work, including a 2026 study of its morphology, but I found no pedagogical materials, standard orthography or dictionary. A curriculum would build on this documentation and on local Konawe Kepulauan schools. Indonesian would be the bridge language.
\end{document}